% \iffalse meta-comment
%
% Copyright (C) SJTUG
%   2018--2026 Weijian Wu   <alexarawu@outlook.com>
%   2022--2026 Zilong Li    <logcreative@outlook.com>
%   2024--2026 Boshi Yuan   <nemoyuan2008@outlook.com>
%
% This work may be distributed and/or modified under the
% conditions of the LaTeX Project Public License, either version 1.3c
% of this license or (at your option) any later version.
% The latest version of this license is in
%   https://www.latex-project.org/lppl.txt
% and version 1.3c or later is part of all distributions of LaTeX
% version 2008 or later.
%
%<*!(driver|install)>
%<+!driver>\GetIdInfo $Id: sjtutex-font.dtx 2.3.1 2026-04-21 15:09:02 +0800 f6eb9dc Alexara Wu <alexarawu@outlook.com>$
%<text&newtx>  {New TX text fonts definition (SJTUTeX)}
%<text&newtx>\ProvidesExplFile{sjtu-text-font-newtx.def}
%<text&newpx>  {New PX text fonts definition (SJTUTeX)}
%<text&newpx>\ProvidesExplFile{sjtu-text-font-newpx.def}
%<text&lm>  {Latin Modern text fonts definition (SJTUTeX)}
%<text&lm>\ProvidesExplFile{sjtu-text-font-lm.def}
%<text&times>  {Times text fonts definition (SJTUTeX)}
%<text&times>\ProvidesExplFile{sjtu-text-font-times.def}
%<text&stixtwo>  {STIX Two text fonts definition (SJTUTeX)}
%<text&stixtwo>\ProvidesExplFile{sjtu-text-font-stixtwo.def}
%<text&libertinus>  {Libertinus text fonts definition (SJTUTeX)}
%<text&libertinus>\ProvidesExplFile{sjtu-text-font-libertinus.def}
%<text&cambria>  {Cambria text fonts definition (SJTUTeX)}
%<text&cambria>\ProvidesExplFile{sjtu-text-font-cambria.def}
%<text&newcm>  {New Computer Modern text fonts definition (SJTUTeX)}
%<text&newcm>\ProvidesExplFile{sjtu-text-font-newcm.def}
%<text&xits>  {XITS text fonts definition (SJTUTeX)}
%<text&xits>\ProvidesExplFile{sjtu-text-font-xits.def}
%<math&newtx>  {New TX math fonts definition (SJTUTeX)}
%<math&newtx>\ProvidesExplFile{sjtu-math-font-newtx.def}
%<math&newpx>  {New PX math fonts definition (SJTUTeX)}
%<math&newpx>\ProvidesExplFile{sjtu-math-font-newpx.def}
%<math&lm>  {Latin Modern math fonts definition (SJTUTeX)}
%<math&lm>\ProvidesExplFile{sjtu-math-font-lm.def}
%<math&times>  {Times math fonts definition (SJTUTeX)}
%<math&times>\ProvidesExplFile{sjtu-math-font-times.def}
%<math&stixtwo>  {STIX Two math fonts definition (SJTUTeX)}
%<math&stixtwo>\ProvidesExplFile{sjtu-math-font-stixtwo.def}
%<math&libertinus>  {Libertinus math fonts definition (SJTUTeX)}
%<math&libertinus>\ProvidesExplFile{sjtu-math-font-libertinus.def}
%<math&cambria>  {Cambria math fonts definition (SJTUTeX)}
%<math&cambria>\ProvidesExplFile{sjtu-math-font-cambria.def}
%<math&newcm>  {New Computer Modern math fonts definition (SJTUTeX)}
%<math&newcm>\ProvidesExplFile{sjtu-math-font-newcm.def}
%<math&xits>  {XITS math fonts definition (SJTUTeX)}
%<math&xits>\ProvidesExplFile{sjtu-math-font-xits.def}
%<cjk&windows>  {Windows CJK fonts definition (SJTUTeX)}
%<cjk&windows>\ProvidesExplFile{sjtu-cjk-font-windows.def}
%<cjk&mac>  {macOS CJK fonts definition (SJTUTeX)}
%<cjk&mac>\ProvidesExplFile{sjtu-cjk-font-mac.def}
%<cjk&ubuntu>  {Ubuntu CJK fonts definition (SJTUTeX)}
%<cjk&ubuntu>\ProvidesExplFile{sjtu-cjk-font-ubuntu.def}
%<cjk&adobe>  {Adobe CJK fonts definition (SJTUTeX)}
%<cjk&adobe>\ProvidesExplFile{sjtu-cjk-font-adobe.def}
%<cjk&fandol>  {Fandol CJK fonts definition (SJTUTeX)}
%<cjk&fandol>\ProvidesExplFile{sjtu-cjk-font-fandol.def}
%<cjk&founder>  {Founder CJK fonts definition (SJTUTeX)}
%<cjk&founder>\ProvidesExplFile{sjtu-cjk-font-founder.def}
%<cjk&hanyi>  {Hanyi CJK fonts definition (SJTUTeX)}
%<cjk&hanyi>\ProvidesExplFile{sjtu-cjk-font-hanyi.def}
%<!driver>  {\ExplFileDate}{\ExplFileVersion}{\ExplFileDescription}
%</!(driver|install)>
%
% \fi
%
% \begin{implementation}
%
% \paragraph*{\fbox{\file{sjtutex-font.dtx}}}
%
% 本文件定义各字库组合下的字体配置。字体配置文件在文档类载入时由
% \tn{load_font} 按需读入，其实现见 \file{sjtutex.dtx} 中的字体配置小节。
%
% \subsection{预定义字体}
%
% \subsubsection{西文与数学字体}
%
% \changes{v2.0.3}{2023/09/25}{新增 \opt{libertinus} 字体配置。}
%    \begin{macrocode}
%<*math|text>
%<@@=sjtu>
%<*math&type1>
%<*newtx|newpx>
\tl_set_eq:NN \l_@@_save_encodingdefault_tl \encodingdefault
\tl_set_eq:NN \l_@@_save_rmdefault_tl \rmdefault
\tl_set_eq:NN \l_@@_save_sfdefault_tl \sfdefault
\tl_set_eq:NN \l_@@_save_ttdefault_tl \ttdefault
\tl_set:Nn \encodingdefault { OT1 }
%<newtx>\tl_set:Nn \rmdefault { ntxtlf }
%<newpx>\tl_set:Nn \rmdefault { zplTLF }
\tl_set:Nn \qhv@scale { 0.94 }
\tl_set:Nn \sfdefault { qhv }
\tl_set:Nn \ttdefault { qcr }
%<newtx>\RequirePackage { newtxmath }
%<newpx>\RequirePackage { newpxmath }
\tl_set_eq:NN \encodingdefault \l_@@_save_encodingdefault_tl
\tl_set_eq:NN \rmdefault \l_@@_save_rmdefault_tl
\tl_set_eq:NN \sfdefault \l_@@_save_sfdefault_tl
\tl_set_eq:NN \ttdefault \l_@@_save_ttdefault_tl
%</newtx|newpx>
%<times>\PassOptionsToPackage { Symbol } { upgreek }
%<lm|times>\RequirePackage { amssymb, upgreek }
%<*lm>
\SetSymbolFont { operators    } { normal } { OT1 } { lmr  } { m  } { n  }
\SetSymbolFont { letters      } { normal } { OML } { lmm  } { m  } { it }
\SetSymbolFont { symbols      } { normal } { OMS } { lmsy } { m  } { n  }
\SetSymbolFont { largesymbols } { normal } { OMX } { lmex } { m  } { n  }
\SetSymbolFont { operators    } { bold   } { OT1 } { lmr  } { bx } { n  }
\SetSymbolFont { letters      } { bold   } { OML } { lmm  } { b  } { it }
\SetSymbolFont { symbols      } { bold   } { OMS } { lmsy } { b  } { n  }
\SetSymbolFont { largesymbols } { bold   } { OMX } { lmex } { m  } { n  }
\SetMathAlphabet { \mathbf } { normal } { OT1 } { lmr  } { bx } { n  }
\SetMathAlphabet { \mathsf } { normal } { OT1 } { lmss } { m  } { n  }
\SetMathAlphabet { \mathit } { normal } { OT1 } { lmr  } { m  } { it }
\SetMathAlphabet { \mathtt } { normal } { OT1 } { lmtt } { m  } { n  }
\SetMathAlphabet { \mathbf } { bold   } { OT1 } { lmr  } { bx } { n  }
\SetMathAlphabet { \mathsf } { bold   } { OT1 } { lmss } { bx } { n  }
\SetMathAlphabet { \mathit } { bold   } { OT1 } { lmr  } { bx } { it }
\SetMathAlphabet { \mathtt } { bold   } { OT1 } { lmtt } { m  } { n  }
\@@_set_slanted_greek:
%</lm>
%<*times>
\tl_set_eq:NN \l_@@_save_rmdefault_tl \rmdefault
  \RequirePackage { mathptmx }
\tl_set_eq:NN \rmdefault \l_@@_save_rmdefault_tl
\tl_set:Nn \Hv@scale { 0.94 }
\DeclareMathAlphabet { \mathsf } { OT1 } { phv } { m } { n }
\DeclareMathAlphabet { \mathtt } { OT1 } { pcr } { m } { n }
\SetMathAlphabet { \mathsf } { bold } { OT1 } { phv } { b } { n }
\SetMathAlphabet { \mathtt } { bold } { OT1 } { pcr } { b } { n }
\DeclareSymbolFont { SJTU@ptm } { OML } { ptmcm } { m } { it }
\@@_declare_math_symbol:nnNn { \mathord } { SJTU@ptm } \upvarsigma { "26 }
%</times>
%<lm|times>\bool_if:NT \g_@@_upright_integral_bool
%<lm|times>  { \RequirePackage { cmupint } }
\@@_set_unimath_symbol:
%</math&type1>
%<*!(math&type1)>
%<lm>\@@_if_engine_opentype:F
%<*!lm>
\@@_if_engine_opentype:TF
  {
%<*math>
    \RequirePackage { unicode-math }
%<*stixtwo>
    \bool_if:NTF \g_@@_upright_integral_bool
      {
        \setmathfont { STIXTwoMath-Regular.otf }
          [ StylisticSet = 8 ]
      }
      { \setmathfont { STIXTwoMath-Regular.otf } }
    \setmathfont { STIXTwoMath-Regular.otf }
      [
        range        = { scr, bfscr },
        StylisticSet = 1
      ]
%</stixtwo>
%<*libertinus>
    \bool_if:NTF \g_@@_upright_integral_bool
      { \setmathfont { LibertinusMath-Regular.otf } }
      {
        \setmathfont { LibertinusMath-Regular.otf }
          [ StylisticSet = 8 ]
      }
    \setmathfont { latinmodern-math.otf } [ range = \checkmark ]
%</libertinus>
%<cambria>    \setmathfont { Cambria~Math }
%<*newcm>
    \bool_if:NTF \g_@@_upright_integral_bool
      {
        \setmathfont { NewCMMath-Book.otf }
          [ StylisticSet = 2 ]
      }
      { \setmathfont { NewCMMath-Book.otf } }
    \setmathfont { NewCMMath-Book.otf }
      [
        range        = { scr, bfscr },
        StylisticSet = 1
      ]
%</newcm>
%<*xits>
    \bool_if:NTF \g_@@_upright_integral_bool
      {
        \setmathfont { XITSMath-Regular }
          [
            Extension    = .otf,
            BoldFont     = XITSMath-Bold,
            StylisticSet = 8
          ]
      }
      {
        \setmathfont { XITSMath-Regular }
          [
            Extension    = .otf,
            BoldFont     = XITSMath-Bold,
          ]
      }
    \setmathfont { XITSMath-Regular.otf }
      [
        range        = { cal, bfcal },
        StylisticSet = 1
      ]
%</xits>
%</math>
%<*newtx|newpx|stixtwo|xits>
%<math>    \setmathrm
%<text>    \setmainfont
%<newtx>      { TeXGyreTermesX }
%<newpx>      { TeXGyrePagellaX }
%<stixtwo>      { STIXTwoText }
%<xits>      { XITS }
      [
        Extension      = .otf,
        UprightFont    = *-Regular,
        BoldFont       = *-Bold,
        ItalicFont     = *-Italic,
        BoldItalicFont = *-BoldItalic
      ]
%<math>    \setmathsf
%<text>    \setsansfont
      { texgyreheros }
      [
        Extension      = .otf,
        UprightFont    = *-regular,
        BoldFont       = *-bold,
        ItalicFont     = *-italic,
        BoldItalicFont = *-bolditalic,
        Scale          = 0.94
      ]
%<math>    \setmathtt
%<text>    \setmonofont
      { texgyrecursor }
      [
        Extension      = .otf,
        UprightFont    = *-regular,
        BoldFont       = *-bold,
        ItalicFont     = *-italic,
        BoldItalicFont = *-bolditalic,
        Ligatures      = CommonOff
      ]
%</newtx|newpx|stixtwo|xits>
%<*text&times>
    \setmainfont { Times~New~Roman } [ Ligatures = Rare ]
    \setsansfont { Arial } [ Scale = 0.94 ]
    \setmonofont { Courier~New }
%</text&times>
%<*libertinus>
%<math>    \setmathrm
%<text>    \setmainfont
      { LibertinusSerif }
      [
        Extension           = .otf,
        UprightFont         = *-Regular,
        BoldFont            = *-Bold,
        ItalicFont          = *-Italic,
        BoldItalicFont      = *-BoldItalic,
        SlantedFont         = *-Regular,
        BoldSlantedFont     = *-Bold,
        SlantedFeatures     = { FakeSlant = 0.2 },
        BoldSlantedFeatures = { FakeSlant = 0.2 }
      ]
%<math>    \setmathsf
%<text>    \setsansfont
      { LibertinusSans }
      [
        Extension           = .otf,
        UprightFont         = *-Regular,
        BoldFont            = *-Bold,
        ItalicFont          = *-Italic,
        BoldItalicFont      = *-Italic,
        BoldItalicFeatures  = { FakeBold  = 3   },
        SlantedFont         = *-Regular,
        BoldSlantedFont     = *-Bold,
        SlantedFeatures     = { FakeSlant = 0.2 },
        BoldSlantedFeatures = { FakeSlant = 0.2 }
      ]
%</libertinus>
%<*cambria>
%<*math>
    \setmathrm { Cambria }
    \setmathsf { Calibri }
    \setmathtt { Consolas } [ Scale = 0.95 ]
%</math>
%<*text>
    \setmainfont { Cambria }
    \setsansfont { Calibri }
    \setmonofont { Consolas } [ Scale = 0.95 ]
%</text>
%</cambria>
%<*newcm>
%<math>    \setmathrm
%<text>    \setmainfont
      { NewCM10 }
      [
        Extension    = .otf,
        SizeFeatures =
          {
            {
              Size        = -9,
              Font        = NewCM08-Book,
              ItalicFont  = NewCM08-BookItalic,
              SlantedFont = NewCM08-Book,
            },
            { Size        = 9- }
          },
        UprightFont         = *-Book,
        BoldFont            = *-Bold,
        ItalicFont          = *-BookItalic,
        BoldItalicFont      = *-BoldItalic,
        SlantedFont         = *-Book,
        SlantedFeatures     = { FakeSlant = 0.25 },
        BoldSlantedFont     = *-Bold,
        BoldSlantedFeatures = { FakeSlant = 0.25 }
      ]
%<math>    \setmathsf
%<text>    \setsansfont
      { NewCMSans10 }
      [
        Extension    = .otf,
        SizeFeatures =
          {
            {
              Size       = -9,
              Font       = NewCMSans08-Book,
              ItalicFont = NewCMSans08-BookOblique,
            },
            { Size       = 9- }
          },
        UprightFont    = *-Book,
        BoldFont       = *-Bold,
        ItalicFont     = *-BookOblique,
        BoldItalicFont = *-BoldOblique
      ]
%<math>    \setmathtt
%<text>    \setmonofont
      { NewCMMono10 }
      [
        Extension           = .otf,
        UprightFont         = *-Book,
        BoldFont            = *-Bold,
        ItalicFont          = *-BookItalic,
        BoldItalicFont      = *-BoldOblique,
        SlantedFont         = *-Book,
        SlantedFeatures     = { FakeSlant = 0.25 },
        BoldSlantedFont     = *-Bold,
        BoldSlantedFeatures = { FakeSlant = 0.25 }
      ]
%</newcm>
  }
%</!lm>
%<*text&!otf>
  {
    \tl_set:Nn \encodingdefault { T1 }
%<lm>    \tl_set:Nn \rmdefault { lmr  }
%<lm>    \tl_set:Nn \sfdefault { lmss }
%<libertinus>    \tl_set:Nn \rmdefault { LibertinusSerif-TLF }
%<libertinus>    \tl_set:Nn \sfdefault { LibertinusSans-TLF  }
%<lm|libertinus>    \tl_set:Nn \ttdefault { lmtt }
%<newtx>    \PassOptionsToPackage { nohelv, nott } { newtxtext }
%<newpx>    \PassOptionsToPackage { nohelv, nott } { newpxtext }
%<newtx>    \RequirePackage { newtxtext }
%<newpx>    \RequirePackage { newpxtext }
%<*stixtwo>
    \DeclareEncodingSubset { TS1 } { ? } { 0 }
    \UndeclareTextCommand { \textpertenthousand } { T1 }
    \DeclareTextSymbolDefault { \textpertenthousand } { TS1 }
    \tl_set:Nn \rmdefault { stix2 }
%</stixtwo>
%<newtx|newpx|stixtwo>    \tl_set:Nn \qhv@scale { 0.94 }
%<newtx|newpx|stixtwo>    \tl_set:Nn \sfdefault { qhv }
%<newtx|newpx|stixtwo>    \tl_set:Nn \ttdefault { qcr }
%<*times>
    \tl_set:Nn \rmdefault { ptm }
    \tl_set:Nn \Hv@scale { 0.94 }
    \tl_set:Nn \sfdefault { phv }
    \tl_set:Nn \ttdefault { pcr }
%</times>
  }
%</text&!otf>
%<*text&otf>
%<cambria>  { \@@_fontset_error:nn { text } { cambria } }
%<newcm>  { \@@_fontset_error:nn { text } { newcm } }
%<xits>  { \@@_fontset_error:nn { text } { xits } }
%</text&otf>
%<*math&!otf>
  {
%<*libertinus>
    \exp_args:No \PassOptionsToPackage
      { \g_@@_math_font_options_clist } { libertinust1math }
    \RequirePackage { libertinust1math }
%</libertinus>
%<*stixtwo>
    \DeclareSizeFunction { sub } { \sub@sfcnt \@font@info }
    \PassOptionsToPackage { notext } { stix2 }
    \RequirePackage { stix2 }
    \clist_map_inline:nn
      {
        \upalpha      { "0B } ,
        \upbeta       { "0C } ,
        \upgamma      { "0D } ,
        \updelta      { "0E } ,
        \upepsilon    { "0F } ,
        \upzeta       { "10 } ,
        \upeta        { "11 } ,
        \uptheta      { "12 } ,
        \upiota       { "13 } ,
        \upkappa      { "14 } ,
        \uplambda     { "15 } ,
        \upmu         { "16 } ,
        \upnu         { "17 } ,
        \upxi         { "18 } ,
        \uppi         { "19 } ,
        \uprho        { "1A } ,
        \upsigma      { "1B } ,
        \uptau        { "1C } ,
        \upupsilon    { "1D } ,
        \upphi        { "1E } ,
        \upchi        { "1F } ,
        \uppsi        { "20 } ,
        \upomega      { "21 } ,
        \upvarepsilon { "22 } ,
        \upvartheta   { "23 } ,
        \upvarpi      { "24 } ,
        \upvarrho     { "25 } ,
        \upvarsigma   { "26 } ,
        \upvarphi     { "27 }
      }
      { \@@_declare_math_symbol:nnNn { \stix@lcgc } { operators } #1 }
    \@@_set_slanted_greek:
%</stixtwo>
  }
%</math&!otf>
%<*math&otf>
%<cambria>  { \@@_fontset_error:nn { math } { cambria } }
%<newcm>  { \@@_fontset_error:nn { math } { newcm } }
%<xits>  { \@@_fontset_error:nn { math } { xits } }
%</math&otf>
%</!(math&type1)>
%</math|text>
%    \end{macrocode}
%
% \subsubsection{CJK 字体}
%
% \changes{v2.2}{2024/11/28}{新增 \opt{hanyi} 字体配置。}
% 在字体未提供对应粗体的情况下，允许使用伪粗。
%    \begin{macrocode}
%<*cjk>
\@@_if_engine_opentype:TF
  {
    \@@_if_main_lang_ja:TF
      {
%<*windows>
        \setCJKmainfont { MS~Mincho } [ AutoFakeBold = 3 ]
        \setCJKsansfont { MS~Gothic } [ AutoFakeBold = 3 ]
        \setCJKmonofont { MS~Mincho }
        \setCJKfamilyfont { jamin  } { MS~Mincho } [ AutoFakeBold = 3 ]
        \setCJKfamilyfont { jagoth } { MS~Gothic } [ AutoFakeBold = 3 ]
%</windows>
%<*mac>
        \setCJKmainfont { HiraMinProN  }
          [
            UprightFont    = *-W3 ,
            BoldFont       = *-W6
          ]
        \setCJKsansfont { HiraKakuProN }
          [
            UprightFont    = *-W3 ,
            BoldFont       = *-W6
          ]
        \setCJKmonofont { HiraMinProN-W3 }
        \setCJKfamilyfont { jamin  } { HiraMinProN  }
          [
            UprightFont    = *-W3 ,
            BoldFont       = *-W6
          ]
        \setCJKfamilyfont { jagoth } { HiraKakuProN }
          [
            UprightFont    = *-W3 ,
            BoldFont       = *-W6
          ]
%</mac>
%<*ubuntu>
        \setCJKmainfont { Noto~Serif~CJK~JP }
          [
            UprightFont = *~Light ,
            BoldFont    = *~Bold
          ]
        \setCJKsansfont { Noto~Sans~CJK~JP  }
          [
            UprightFont = *~Medium ,
            BoldFont    = *~Bold
          ]
        \setCJKmonofont { Noto~Serif~CJK~JP }
          [
            UprightFont = *~Light ,
            BoldFont    = *~Bold
          ]
        \setCJKfamilyfont { jamin  } { Noto~Serif~CJK~JP }
          [
            UprightFont = *~Light ,
            BoldFont    = *~Bold
          ]
        \setCJKfamilyfont { jagoth } { Noto~Sans~CJK~JP  }
          [
            UprightFont = *~Medium ,
            BoldFont    = *~Bold
          ]
%</ubuntu>
%<*adobe>
        \setCJKmainfont { KozMinPr6N }
          [
            UprightFont = *-Light ,
            BoldFont    = *-Bold
          ]
        \setCJKsansfont { KozGoPr6N  }
          [
            UprightFont = *-Medium ,
            BoldFont    = *-Bold
          ]
        \setCJKmonofont { KozMinPr6N-Light }
        \setCJKfamilyfont { jamin  } { KozMinPr6N }
          [
            UprightFont = *-Light ,
            BoldFont    = *-Bold
          ]
        \setCJKfamilyfont { jagoth } { KozGoPr6N  }
          [
            UprightFont = *-Medium ,
            BoldFont    = *-Bold
          ]
%</adobe>
%<*fandol|hanyi>
        \setCJKmainfont { HaranoAjiMincho }
          [
            Extension   = .otf ,
            UprightFont = *-Regular ,
%<fandol>            BoldFont    = *-Bold
%<hanyi>            BoldFont    = *-Medium
          ]
        \setCJKsansfont { HaranoAjiGothic }
          [
            Extension   = .otf ,
            UprightFont = *-Medium ,
            BoldFont    = *-Bold
          ]
        \setCJKmonofont { HaranoAjiGothic }
          [
            Extension   = .otf ,
            UprightFont = *-Regular
          ]
        \setCJKfamilyfont { jamin  } { HaranoAjiMincho }
          [
            Extension   = .otf ,
%<fandol>            BoldFont    = *-Bold
%<hanyi>            BoldFont    = *-Medium
          ]
        \setCJKfamilyfont { jagoth } { HaranoAjiGothic }
          [
            Extension   = .otf ,
            UprightFont = *-Medium ,
            BoldFont    = *-Bold
          ]
%</fandol|hanyi>
%<*founder>
        \setCJKmainfont { ipam.ttf } [ AutoFakeBold = 3 ]
        \setCJKsansfont { ipag.ttf } [ AutoFakeBold = 3 ]
        \setCJKmonofont { ipag.ttf }
        \setCJKfamilyfont { jamin  } { ipam.ttf } [ AutoFakeBold = 3 ]
        \setCJKfamilyfont { jagoth } { ipag.ttf } [ AutoFakeBold = 3 ]
%</founder>
        \NewDocumentCommand \mincho   { } { \CJKfamily { jamin   } }
        \NewDocumentCommand \gothic   { } { \CJKfamily { jagoth  } }
      }
%<*windows>
      {
        \setCJKmainfont { SimSun   }
          [ AutoFakeBold = 3 , ItalicFont = KaiTi ]
        \setCJKsansfont { SimHei   } [ AutoFakeBold = 3 ]
        \setCJKmonofont { FangSong }
      }
    \setCJKfamilyfont { zhsong } { SimSun   }
      [ AutoFakeBold = 3 , ItalicFont = KaiTi ]
    \setCJKfamilyfont { zhhei  } { SimHei   } [ AutoFakeBold = 3 ]
    \setCJKfamilyfont { zhkai  } { KaiTi    }
    \setCJKfamilyfont { zhfs   } { FangSong }
%</windows>
%<*mac>
      {
        \setCJKmainfont { Songti~SC  }
          [
            UprightFont    = *~Light ,
            BoldFont       = *~Bold ,
            ItalicFont     = Kaiti~SC~Regular ,
            BoldItalicFont = Kaiti~SC~Bold
          ]
        \setCJKsansfont { Heiti~SC   }
          [
            UprightFont    = *~Medium ,
            AutoFakeBold   = 3
          ]
        \setCJKmonofont { STFangsong }
      }
    \setCJKfamilyfont { zhsong } { Songti~SC  }
      [
        UprightFont    = *~Light ,
        BoldFont       = *~Bold ,
        ItalicFont     = Kaiti~SC~Regular ,
        BoldItalicFont = Kaiti~SC~Bold
      ]
    \setCJKfamilyfont { zhhei  } { Heiti~SC   }
      [
        UprightFont    = *~Medium ,
        AutoFakeBold   = 3
      ]
    \setCJKfamilyfont { zhfs   } { STFangsong }
    \setCJKfamilyfont { zhkai  } { Kaiti~SC   }
      [
        UprightFont    = *~Regular ,
        BoldFont       = *~Bold
        ]
%</mac>
%<*ubuntu>
      {
        \setCJKmainfont { Noto~Serif~CJK~SC }
          [
            UprightFont = *~Light ,
            BoldFont    = *~Bold ,
            ItalicFont  = AR~PL~KaitiM~GB
          ]
        \setCJKsansfont { Noto~Sans~CJK~SC  }
          [
            UprightFont = *~Medium ,
            BoldFont    = *~Bold
          ]
        \setCJKmonofont { Noto~Serif~CJK~SC }
          [
            UprightFont = *~Light ,
            BoldFont    = *~Bold
          ]
      }
    \setCJKfamilyfont { zhsong } { Noto~Serif~CJK~SC }
      [
        UprightFont = *~Light ,
        BoldFont    = *~Bold ,
        ItalicFont  = AR~PL~KaitiM~GB
      ]
    \setCJKfamilyfont { zhhei  } { Noto~Sans~CJK~SC  }
      [
        UprightFont = *~Medium ,
        BoldFont    = *~Bold
      ]
    \setCJKfamilyfont { zhkai  } { AR~PL~KaitiM~GB   }
%</ubuntu>
%<*adobe>
      {
        \setCJKmainfont { AdobeSongStd-Light       }
          [ AutoFakeBold = 3 , ItalicFont = AdobeKaitiStd-Regular ]
        \setCJKsansfont { AdobeHeitiStd-Regular    } [ AutoFakeBold = 3 ]
        \setCJKmonofont { AdobeFangsongStd-Regular }
      }
    \setCJKfamilyfont { zhsong } { AdobeSongStd-Light       }
      [ AutoFakeBold = 3 , ItalicFont = AdobeKaitiStd-Regular ]
    \setCJKfamilyfont { zhhei  } { AdobeHeitiStd-Regular    } [ AutoFakeBold = 3 ]
    \setCJKfamilyfont { zhfs   } { AdobeFangsongStd-Regular }
    \setCJKfamilyfont { zhkai  } { AdobeKaitiStd-Regular    }
%</adobe>
%<*fandol>
      {
        \setCJKmainfont { FandolSong }
          [
            Extension   = .otf ,
            UprightFont = *-Regular ,
            BoldFont    = *-Bold ,
            ItalicFont  = FandolKai-Regular
          ]
        \setCJKsansfont { FandolHei  }
          [
            Extension   = .otf ,
            UprightFont = *-Regular ,
            BoldFont    = *-Bold
          ]
        \setCJKmonofont { FandolFang }
          [
            Extension   = .otf ,
            UprightFont = *-Regular
          ]
      }
    \setCJKfamilyfont { zhsong } { FandolSong }
      [
        Extension   = .otf ,
        UprightFont = *-Regular ,
        BoldFont    = *-Bold ,
        ItalicFont  = FandolKai-Regular
      ]
    \setCJKfamilyfont { zhhei  } { FandolHei  }
      [
        Extension   = .otf ,
        UprightFont = *-Regular,
        BoldFont    = *-Bold
      ]
    \setCJKfamilyfont { zhfs   } { FandolFang }
      [
        Extension   = .otf ,
        UprightFont = *-Regular
      ]
    \setCJKfamilyfont { zhkai  } { FandolKai  }
      [
        Extension   = .otf ,
        UprightFont = *-Regular
      ]
%</fandol>
%<*founder>
      {
        \setCJKmainfont { FZShuSong-Z01  }
          [ AutoFakeBold = 3 , ItalicFont = FZKai-Z03 ]
        \setCJKsansfont { FZHei-B01      } [ AutoFakeBold = 3 ]
        \setCJKmonofont { FZFangSong-Z02 }
      }
    \setCJKfamilyfont { zhsong } { FZShuSong-Z01  }
      [ AutoFakeBold = 3 , ItalicFont = FZKai-Z03 ]
    \setCJKfamilyfont { zhhei  } { FZHei-B01      } [ AutoFakeBold = 3 ]
    \setCJKfamilyfont { zhkai  } { FZKai-Z03      }
    \setCJKfamilyfont { zhfs   } { FZFangSong-Z02 }
%</founder>
%<*hanyi>
      {
        \setCJKmainfont { HYShuSongEr~S }
          [ BoldFont = HYZhongSong~S, ItalicFont = HYKaiTi~S ]
        \setCJKsansfont { HYZhongHei~S  } [ BoldFont = HYDaHei~S ]
        \setCJKmonofont { HYFangSong~S  }
      }
    \setCJKfamilyfont { zhsong } { HYShuSongEr~S }
      [ BoldFont = HYZhongSong~S ]
    \setCJKfamilyfont { zhhei  } { HYZhongHei~S  }
      [ BoldFont = HYDaHei~S     ]
    \setCJKfamilyfont { zhkai  } { HYKaiTi~S     }
    \setCJKfamilyfont { zhfs   } { HYFangSong~S  }
%</hanyi>
    \NewDocumentCommand \songti   { } { \CJKfamily { zhsong  } }
    \NewDocumentCommand \heiti    { } { \CJKfamily { zhhei   } }
%<!ubuntu>    \NewDocumentCommand \fangsong { } { \CJKfamily { zhfs    } }
    \NewDocumentCommand \kaishu   { } { \CJKfamily { zhkai   } }
  }
  {
    \@@_if_main_lang_ja:TF
%<windows>      { \@@_fontset_error:nn { cjk } { windows } }
%<windows>      { \ctex_file_input:n { ctex-fontset-windows.def } }
%<mac>      { \@@_fontset_error:nn { cjk } { mac } }
%<mac>      { \ctex_file_input:n { ctex-fontset-mac.def } }
%<ubuntu>      { \@@_fontset_error:nn { cjk } { ubuntu } }
%<ubuntu>      { \ctex_file_input:n { ctex-fontset-ubuntu.def } }
%<adobe>      { \@@_fontset_error:nn { cjk } { adobe } }
%<adobe>      { \ctex_file_input:n { ctex-fontset-adobe.def } }
%<fandol>      { \@@_fontset_error:nn { cjk } { fandol } }
%<fandol>      { \ctex_file_input:n { ctex-fontset-fandol.def } }
%<founder>      { \@@_fontset_error:nn { cjk } { founder } }
%<founder>      { \ctex_file_input:n { ctex-fontset-founder.def } }
%<hanyi>      { \@@_fontset_error:nn { cjk } { hanyi } }
%<hanyi>      { \ctex_file_input:n { ctex-fontset-hanyi.def } }
  }
%</cjk>
%    \end{macrocode}
% \end{implementation}
%
% \endinput
